\section{Pre-Estimation}
\label{sec:pre_estimation}

This section describes the data pipeline from raw observations to estimation-ready objects.
The pipeline handles data loading, validation, state discretization, feature normalization, transition estimation, and utility specification.

\paragraph{Data types.}
Three core types encode the estimation problem.
\texttt{DDCProblem} is a frozen dataclass that stores the MDP structure, including the number of states $|\mathcal{S}|$, the number of actions $|\mathcal{A}|$, the discount factor $\beta$, and the scale parameter $\sigma$ of the extreme value shocks.
\texttt{Trajectory} stores a single individual's sequence of state-action-next-state tuples $(s_t, a_t, s_{t+1})$ as JAX arrays with matching lengths.
\texttt{Panel} collects a list of trajectories and provides cached sufficient statistics via \texttt{.sufficient\_stats()}, which computes state-action counts, empirical CCPs, transition frequencies, and the initial state distribution in a single pass over the data.
Panels can be constructed from pandas DataFrames via \texttt{Panel.from\_dataframe(df, state, action, id, next\_state)} or saved and loaded as compressed NumPy archives via \texttt{.save\_npz()} and \texttt{.load\_npz()} for reproducibility.

\paragraph{Data loading.}
The \texttt{econirl.datasets} module provides loaders for real-world datasets spanning DDC applications and IRL trajectory data.
DDC datasets include the Rust (1987) bus engine replacement data, Scania heavy truck component reliability, Keane-Wolpin career choices, and ICU sepsis treatment protocols.
IRL trajectory datasets include Beijing taxi GPS traces via the T-Drive corpus, NYC Citibike station-to-station route choices, NGSIM highway lane-change trajectories, and Foursquare check-in sequences.
Each loader returns a Panel object with metadata documenting the source, sample period, and variable definitions.
For synthetic data, \texttt{simulate\_panel(env, n\_individuals, n\_periods, seed)} generates panels from any \texttt{DDCEnvironment} with known ground truth parameters, enabling controlled parameter recovery experiments.

\paragraph{Data validation.}
Before estimation, \texttt{check\_panel\_structure()} validates the panel for missing values, unbalanced periods, and period gaps, returning a \texttt{PanelValidationResult} with error and warning lists.
Four pre-estimation diagnostics are printed automatically before any \texttt{.estimate()} call.
Feature matrix rank is computed on the flattened feature matrix of dimension $(|\mathcal{S}| \times |\mathcal{A}|) \times K$, and if the rank falls below $K$, estimation halts with an error.
The condition number flags near-collinearity that inflates standard errors even when the matrix is technically full rank, with values above $10^6$ triggering a warning.
State coverage reports how many states have at least one observation.
Single-action state counts identify states where only one action is ever observed, which produce degenerate CCP estimates.

\paragraph{State discretization and normalization.}
Continuous state variables are discretized via \texttt{discretize\_state(values, method, n\_bins)} using either equal-width or equal-count binning.
The Rust (1987) convention of 5{,}000-mile bins is available via \texttt{discretize\_mileage()}.
Feature normalization during training uses \texttt{RunningNorm}, which implements the numerically stable online mean-variance algorithm of Chan, Golub, and LeVeque (1979) for features that span different scales.
This is important when feature columns differ by orders of magnitude, as in the Rust bus model where operating cost is $O(10^{-3})$ and replacement cost is $O(1)$.

\paragraph{Transition estimation.}
When transition probabilities are not known from the model, they are estimated from data via \texttt{estimate\_transition\_probs()}, which counts state transitions from non-replacement observations and returns a probability vector that is expanded into the full $(|\mathcal{A}|, |\mathcal{S}|, |\mathcal{S}|)$ transition tensor.
Group-stratified estimation is available via \texttt{estimate\_transition\_probs\_by\_group()} for heterogeneous transition dynamics across subpopulations such as bus groups or geographic regions.
All estimators expect transitions in the shape $(|\mathcal{A}|, |\mathcal{S}|, |\mathcal{S}|)$, meaning \texttt{transitions[a, s, s']} gives $P(s'|s,a)$.

\paragraph{Utility specification.}
The \texttt{LinearUtility} class specifies $u(s,a;\theta) = \theta^\top \phi(s,a)$ from a feature matrix of shape $(|\mathcal{S}|, |\mathcal{A}|, K)$.
The \texttt{anchor\_action} parameter normalizes utility differences by subtracting the anchor action's features from all other actions, resolving the level non-identification discussed in Section~\ref{sec:framework}.
The class method \texttt{LinearUtility.from\_environment(env)} extracts features and parameter names automatically from any \texttt{DDCEnvironment}.
For deep IRL estimators that learn nonlinear reward functions, \texttt{NeuralCost} provides a neural network parameterization of the reward.
